\documentclass[11pt]{article}
\usepackage[a4paper,margin=27mm]{geometry}
\usepackage{amsmath,amssymb,amsthm,mathtools,microtype}
\usepackage[colorlinks=true,allcolors=blue]{hyperref}
\usepackage{booktabs,enumitem}
\setlist{nosep,leftmargin=*,topsep=2pt}
\newtheorem{theorem}{Theorem}[section]
\newtheorem{proposition}[theorem]{Proposition}
\newtheorem{lemma}[theorem]{Lemma}
\newtheorem{corollary}[theorem]{Corollary}
\theoremstyle{remark}\newtheorem{remark}[theorem]{Remark}
\newcommand{\Herm}{\operatorname{Herm}}
\newcommand{\Tr}{\operatorname{Tr}}
\newcommand{\id}{\mathbf 1}
\newcommand{\so}{\mathfrak{so}}
\title{Positive Cartan Carriers and Signed Lorentz Selectors\\[0.4em]\large Minimal spinors and exact restricted-action certificates}
\author{}
\date{29 October 2026}

\hypersetup{pdfauthor={},pdftitle={Positive Cartan Carriers and Signed Lorentz Selectors}}
\begin{document}
\maketitle

\begin{abstract}
A positive carrier for a symmetry need not contain a positive order parameter
with the desired stabilizer.
We examine this distinction for the de Sitter and anti-de Sitter Cartan
algebras in four spacetime dimensions.
One qubit cone admits no nontrivial reversible linear action of either
simple ten-generator algebra, while the classical four-component spinor is
minimal among faithful complex-linear carriers.
Its congruence action preserves the positive cone, but its Lorentz contact
orbit is signed.
More strongly, every nonzero Lorentz-fixed Hermitian matrix in the stated
Dirac restriction has inertia $(2,2)$, excluding a positive semidefinite
selector even on the boundary.
We give explicit Clifford matrices, analytic proofs, and portable exact
certificates for the one-dimensional contact-linear curvature-square
invariant, the independent topological quadratic direction, torsion
invertibility and the Einstein/Weyl rank split.
The Clifford carrier and restricted Cartan action are established ingredients,
not new gravity actions; their role here is to locate precisely the
positive-carrier/selector obstruction and its conditional classical bridge.
No contact selection law, coupling scale, global Spin lift or quantum measure
is derived.
\end{abstract}

\section{Question, conventions, and claim boundary}

MacDowell--Mansouri and Stelle--West gravity use a de Sitter or anti-de Sitter gauge
connection together with a reduction to the Lorentz subgroup
\cite{MacDowellMansouri,StelleWest,Wise}.  Modern Cartan-gravity formulations make the
contact vector explicit and may promote it to a dynamical field
\cite{WestmanZlosnikMatter,WestmanZlosnikExplore,WestmanZlosnikIntro}.  Here we study a
more restricted algebraic question: can one finite positive carrier realize the full
Cartan group and also select its Lorentzian gravitational orbit?

Let
\begin{equation}
 \eta_{AB}=\operatorname{diag}(1,-1,-1,-1,s),\qquad
 (J_{AB})^C{}_D=\delta_A^C\eta_{BD}-\delta_B^C\eta_{AD},
 \qquad P_I=\ell^{-1}J_{I4},
 \label{eq:cartan-generators}
\end{equation}
where $s=\pm1$, $I=0,\ldots,3$, and $\ell>0$.  Direct multiplication gives
\begin{equation}
 [J_{IJ},P_K]=\eta_{JK}P_I-\eta_{IK}P_J,\qquad
 [P_I,P_J]=-\frac{s}{\ell^2}J_{IJ}.
 \label{eq:cartan-brackets}
\end{equation}
Thus $s=-1$ gives $\so(1,4)$ and $s=+1$ gives $\so(2,3)$ in the displayed internal
signature.  Physical de Sitter/anti-de Sitter names also depend on the spacetime metric
sign convention; we do not infer them from $s$ alone.

The main theorems are finite and local.  ``Unique'' below always means within the explicitly
declared class: oriented four-forms, quadratic in the full Cartan curvature, at most linear
in the contact field, and containing neither $DV$ nor higher derivatives.  The theorem does
not classify nonlocal, higher-curvature, fermionic, or boundary actions.

\section{The Cartan deformation is one-parameter}

Let $V$ be the real Lorentz vector representation and
$\mathfrak h=\so(1,3)\simeq\Lambda^2V$.  A Lorentz-equivariant antisymmetric translation
bracket is a map $\Lambda^2V\to\mathfrak h$.  Because the Lorentzian Hodge operator on real
bivectors satisfies $\star^2=-1$, the real commutant is two-dimensional and the most general
candidate is
\begin{equation}
 [P_I,P_J]=\lambda J_{IJ}+\mu(\star J)_{IJ}.
 \label{eq:dual-bracket}
\end{equation}

\begin{proposition}[Jacobi removes the dual deformation]
\label{prop:jacobi}
The three-translation Jacobi identity for \eqref{eq:dual-bracket} implies $\mu=0$.
Consequently the nonzero deformation is, up to scale and real form, the standard
de Sitter/anti-de Sitter bracket.
\end{proposition}

\begin{proof}
The ordinary $\lambda$ terms cancel cyclically, while the dual terms give
\begin{equation}
 \sum_{\rm cyc}[[P_I,P_J],P_K]=3\mu\,\epsilon_{IJKL}P^L.
 \label{eq:jacobi-fault}
\end{equation}
Linear independence of the $P_L$ forces $\mu=0$.
\end{proof}

This preliminary result prevents a spurious parity-odd local translation bracket from
being counted as a second Cartan algebra.  For either simple real form the vector
representation has $B(X,Y)=3\Tr_5(XY)$ and
\begin{equation}
 B(P_I,P_J)=-\frac{6s}{\ell^2}\eta_{IJ}.
 \label{eq:killing-form}
\end{equation}
The $\ell\to\infty$ In\"on\"u--Wigner contraction removes both $[P,P]$ and the translation
block of the Killing form as $\ell^{-2}$.  This is an algebra contraction, not a spacetime
continuum limit.

\section{No full simple Cartan action on one qubit cone}

The connected automorphism algebra of the qubit positive cone has dimension
seven.
A direct proof avoids assuming a classification for larger matrix cones.
Identify the qubit cone with $x^0\geq|\boldsymbol x|$ and set
$q(x)=(x^0)^2-|\boldsymbol x|^2$.
An invertible real cone automorphism $T$ preserves its boundary, so the
quadratic polynomial $q(Tx)$ vanishes on every $(1,\boldsymbol n)$ with
$|\boldsymbol n|=1$.
Its linear term on the sphere vanishes and its spatial quadratic part is a
scalar multiple of $|\boldsymbol n|^2$, giving $q(Tx)=a q(x)$.
Interior preservation implies $a>0$.
In the identity component $T=\sqrt a\,\Lambda$ with
$\Lambda\in SO^+(1,3)$, hence dimension $1+6=7$.
Equivalently, the congruence action of $GL(2,\mathbb C)$ has an imaginary
scalar kernel and the same effective dimension; its determinant-preserving
part is the standard Lorentz spin action~\cite{ArrighiPatricotLorentz}.

\begin{theorem}[One-cone obstruction]
\label{thm:qubit-obstruction}
There is no nontrivial reversible linear action of $\so(1,4)$ or $\so(2,3)$ on one
$\Herm(2)_+$ cone.
\end{theorem}

\begin{proof}
The kernel of a homomorphism from either ten-dimensional simple algebra is an ideal.  A
nonzero homomorphism into the seven-dimensional effective cone-automorphism algebra would
therefore be injective, which is impossible by dimension.
\end{proof}

The theorem explains a structural split: one qubit cone carries the Lorentz subgroup, but
the four coset directions cannot all be reversible positive transformations of the same
qubit.  They may be signed response directions, or the carrier must be enlarged.  The
statement does not exclude nonlinear, relational, or multi-event actions.

\section{Four complex components are minimal and sufficient}

Any complex-linear representation of either real form extends to the simple complex
algebra $\so(5,\mathbb C)$ of dimension ten.  A nonzero representation is faithful, and its
generators are traceless.  It must therefore embed into $\mathfrak{sl}(n,\mathbb C)$.  Since
\begin{equation}
 \dim_\mathbb C\mathfrak{sl}(2,\mathbb C)=3,\qquad
 \dim_\mathbb C\mathfrak{sl}(3,\mathbb C)=8,
 \label{eq:minimal-dimension}
\end{equation}
$n=2,3$ are impossible.

For an explicit construction, let $\sigma_1,\sigma_2,\sigma_3$ be the Pauli
matrices and define
\begin{align}
 E_0&=\sigma_1\otimes I_2,&E_1&=\sigma_2\otimes I_2,\nonumber\\
 E_2&=\sigma_3\otimes\sigma_1,&E_3&=\sigma_3\otimes\sigma_2,
 &E_4&=\sigma_3\otimes\sigma_3 .
 \label{eq:euclidean-gamma}
\end{align}
They anticommute and square to $I_4$.
Set
\begin{equation}
 (\Gamma_0,\Gamma_1,\Gamma_2,\Gamma_3,\Gamma_4)
 =(E_0,iE_1,iE_2,iE_3,\zeta_s E_4),\qquad
 \zeta_{-1}=i,\quad\zeta_{+1}=1.
 \label{eq:explicit-gamma}
\end{equation}
The resulting $4\times4$ Clifford matrices satisfy

\begin{equation}
 \{\Gamma_A,\Gamma_B\}=2\eta_{AB}\id_4,\qquad
 S_{AB}=\frac14[\Gamma_A,\Gamma_B].
 \label{eq:spin-generators}
\end{equation}
They obey
\begin{equation}
 [S_{AB},S_{CD}]=\eta_{BC}S_{AD}-\eta_{AC}S_{BD}
 -\eta_{BD}S_{AC}+\eta_{AD}S_{BC}.
 \label{eq:spin-bracket}
\end{equation}

\begin{theorem}[Minimal positive carrier]
\label{thm:minimal-carrier}
Four complex components are the smallest faithful complex-linear carrier for either
Cartan real form.  The spin representation \eqref{eq:spin-generators} realizes all ten
generators and induces reversible cone automorphisms
\begin{equation}
 X\longmapsto GXG^\dagger,\qquad X\in\Herm(4)_+.
 \label{eq:four-cone-action}
\end{equation}
The central spin element $-\id_4$ acts trivially on the cone.
\end{theorem}
\begin{proof}
The preceding complexification argument excludes dimensions two and three;
a one-dimensional representation is trivial because the algebra is perfect.
The Clifford anticommutation relations give the displayed spin commutators
by direct multiplication.
The representation is nonzero, so simplicity makes its Lie-algebra kernel
zero; thus its ten generators are independent.
Congruence by every invertible spin matrix preserves positive
semidefiniteness in both directions.
Finally $(-I_4)X(-I_4)^\dagger=X$.
\end{proof}

Four-component Cartan spinors and their Clifford realization are classical;
for example, Ref.~\cite{WestmanZlosnikIntro}, Section~5.5.1, explicitly
contrasts them with the two-component Lorentz representation.
The point here is the compatibility question between that carrier and a
positive selector, not a new discovery of the spin representation.


This is a symmetry of the unnormalised cone or of a moving-order-unit structure.  A generic
noncompact $G$ is not unitary and \eqref{eq:four-cone-action} is not a deterministic
trace-preserving channel on normalized four-level states.

\section{The Lorentz contact orbit is necessarily signed}

The explicit matrices
\begin{equation}
 C_{-1}=\Gamma_0,\qquad C_{+1}=i\Gamma_0\Gamma_4,
 \qquad\alpha_{-1}=1,\quad\alpha_{+1}=i
 \label{eq:contact-conventions}
\end{equation}
satisfy $S_{BC}C_s+C_sS_{BC}^\dagger=0$.
They give Hermitian contact matrices by
\begin{equation}
 Q_A=\alpha_s\Gamma_AC_s=Q_A^\dagger,\qquad Q(V)=V^AQ_A,
 \label{eq:contact-matrices}
\end{equation}
and
\begin{align}
 S_{BC}Q_A+Q_AS_{BC}^\dagger
 &=\eta_{CA}Q_B-\eta_{BA}Q_C,\nonumber\\
 \det Q(V)&=(V^A\eta_{AB}V^B)^2.
 \label{eq:contact-covariance}
\end{align}
At $V=\ell\delta_4$, the six $S_{IJ}$ stabilize $Q(V)$ and the four $S_{I4}$ span the
contact orbit.

\begin{theorem}[Positivity--gravity separation]
\label{thm:signed-contact}
For both real forms, $Q(\ell\delta_4)$ has inertia $(2,2,0)$ and is not positive or
negative.  Conversely the stabilizer under congruence of any $X>0$ is conjugate to a
unitary group and hence compact.  Therefore no interior positive state or order unit can
have the required noncompact $Spin(1,3)$ stabilizer.
\end{theorem}

\begin{proof}
The spectrum of the explicit Clifford carrier gives two positive and two negative
eigenvalues.  For the converse, the condition $GXG^\dagger=X$ is equivalent to
$UU^\dagger=I$ for $U=X^{-1/2}GX^{1/2}$.
Thus the full congruence stabilizer is conjugate to $U(4)$; its closed
subgroups are compact.
\end{proof}

Thus positivity can carry the transformations, but the field selecting Lorentzian Cartan
gravity is a signed Hermitian order parameter.  The theorem is a no-go only for an interior
positive selector in a linear congruence carrier; it does not forbid a relational or
boundary selection law.

\subsection{The boundary of the cone does not evade the selector obstruction}

The preceding compact-stabilizer argument applies to positive-definite
matrices in any finite congruence carrier.
For the particular four-component Dirac carrier, the statement can be
strengthened to include its positive semidefinite boundary.

\begin{theorem}[All Lorentz-fixed Hermitian selectors]
\label{thm:boundary-selector}
In a Weyl basis the connected Lorentz subgroup acts by
\begin{equation}
 G(A)=\operatorname{diag}\bigl(A,(A^\dagger)^{-1}\bigr),
 \qquad A\in SL(2,\mathbb C).
 \label{eq:weyl-action}
\end{equation}
Its fixed Hermitian matrices under $X\mapsto G(A)XG(A)^\dagger$ are exactly
\begin{equation}
 X_c=\begin{pmatrix}0&cI_2\\\overline cI_2&0\end{pmatrix},
 \qquad c\in\mathbb C.
 \label{eq:fixed-selector}
\end{equation}
Every nonzero such matrix has inertia $(2,2,0)$.
In particular, zero is the only Lorentz-fixed positive semidefinite matrix.
\end{theorem}
\begin{proof}
For $A\in SU(2)$ the two blocks in \eqref{eq:weyl-action} coincide.
The defining two-dimensional representation is irreducible, so rotation
invariance forces a Hermitian fixed matrix into the form
\[
 X=\begin{pmatrix}aI_2&cI_2\\\overline cI_2&dI_2\end{pmatrix},
 \qquad a,d\in\mathbb R.
\]
Apply the real boost $A_t=\operatorname{diag}(e^{t/2},e^{-t/2})$.
The diagonal blocks become $aA_t^2$ and $dA_t^{-2}$.
Invariance for every $t$ forces $a=d=0$.
Conversely the upper off-diagonal block transforms as
$A(cI_2)A^{-1}=cI_2$ for every $A\in SL(2,\mathbb C)$, and similarly for
the lower block, so every \eqref{eq:fixed-selector} is fixed.
Finally $X_c^2=|c|^2I_4$ and $\Tr X_c=0$, giving two eigenvalues
$+|c|$ and two eigenvalues $-|c|$ when $c\ne0$.
\end{proof}

The Lorentz restriction of both Cartan spinor real forms is equivalent to
\eqref{eq:weyl-action}; the conclusion is therefore the same for both.
The theorem excludes positive fixed selectors in this carrier, including
singular ones.
It does not exclude a relational selection rule, an enlarged representation
with a trivial summand, a nonlinear action, or independently supplied
boundary data.
A signed order parameter is compatible with a positive state space because
observables and states are different objects.


\section{Clifford trace and restricted action selection}

The carrier also supplies the all-index trace identity
\begin{equation}
 \Tr(\Gamma_ES_{AB}S_{CD})=c_s\epsilon_{EABCD},
 \qquad c_{-1}=-1,\quad c_{+1}=i,
 \label{eq:clifford-trace}
\end{equation}
up to the displayed representation convention.  Hence
\begin{equation}
 S_{\rm MM}[A,V]=\alpha\int\Tr\!\left(\Gamma(V)\,\mathcal F\wedge\mathcal F\right)
 \label{eq:mm-action}
\end{equation}
is the contact-weighted MacDowell--Mansouri density.

\begin{theorem}[Invariant count in the declared class]
\label{thm:invariant-count}
Among oriented local four-forms linear in $V$, quadratic in the ten-component Cartan
curvature, and containing no $DV$, the invariant coefficient space is one-dimensional.
It is spanned by \eqref{eq:clifford-trace}.
\end{theorem}

\begin{proof}
Complexify to $SO(5,\mathbb C)$ and use an orthonormal complex basis.
An invariant rank-five tensor is fixed by every diagonal sign matrix having
an even number of minus signs.
For a nonzero coefficient, the five index multiplicities must therefore
have the same parity: flipping labels $a,b$ gives the condition
$m_a+m_b=0\pmod2$.
Their total is five, so every multiplicity is odd, forcing each of the five
labels to occur exactly once.
Determinant-one signed permutation matrices then force the coefficients to
alternate, leaving only a multiple of $\epsilon_{ABCDE}$.
This tensor is antisymmetric in each bivector pair and symmetric under
exchange of the pairs, so it belongs to
$V\otimes\operatorname{Sym}^2(\Lambda^2V)$ and spans its invariant space.
The real epsilon tensor gives the corresponding one-dimensional real
invariant for both real forms.
\end{proof}

The companion integer constraint has 275 columns and rank 274, verified by
an exact modular lower bound and an integer epsilon null vector as specified
in Appendix~\ref{app:certificates}.
The action-space answer itself is prior content.
The polynomial family in Ref.~\cite{WestmanZlosnikIntro}, equations
(82)--(89), contains the coefficients
$a_1\epsilon_{ABCDE}V^E+a_2V_AV_C\eta_{BD}+a_3\eta_{AC}\eta_{BD}$
in its curvature-square sector.
Restricting to degree at most one in $V$ and no $DV$ leaves exactly the
epsilon-linear term and a constant Killing term.
This theorem and the finite certificate make that restricted answer
explicit; they do not introduce a new gravitational action or derive the
restriction from microscopic physics.

In the contact gauge $V=\ell\delta_4$, write
\begin{equation}
 F^{IJ}=R^{IJ}+\lambda e^I\wedge e^J,\qquad
 F^{I4}=T^I/\ell,\qquad \lambda=-s/\ell^2.
 \label{eq:curvature-split}
\end{equation}
The contact-linear epsilon density contains only $F^{IJ}$ and expands into Euler,
Einstein--Cartan, and cosmological terms with fixed relative coefficients.
Explicitly, after absorbing the nonzero Clifford trace constant into a free
coefficient $\beta$, the reduced action is
\begin{align}
 S_\epsilon
 &=\beta\int\epsilon_{IJKL}F_h^{IJ}\wedge F_h^{KL}\nonumber\\
 &=\beta\int\epsilon_{IJKL}\bigl(
 R^{IJ}\wedge R^{KL}
 +2\lambda e^I\wedge e^J\wedge R^{KL}
 +\lambda^2 e^I\wedge e^J\wedge e^K\wedge e^L\bigr).
 \label{eq:explicit-action}
\end{align}
For $\lambda\ne0$, the choice $\beta=1/(8\kappa\lambda)$ makes the
Einstein--Cartan coefficient $1/(4\kappa)$ and gives the internal-convention
cosmological coefficient $\Lambda=-3\lambda$.
This identifies coefficients once $\kappa$ is supplied; it does not derive
Newton's normalization.
The degree-zero quadratic invariant is unique because the complexified
adjoint representation of the simple algebra $\mathfrak{so}(5,\mathbb C)$
has a one-dimensional invariant bilinear space.

There is also the full-Cartan Killing/Chern--Weil density
\begin{equation}
 P_G=\sum_{A<B}\eta_{AA}\eta_{BB}F^{AB}\wedge F^{AB}.
 \label{eq:chern-weil}
\end{equation}
Define $H=\sum_{I<J}R_{IJ}\wedge e^I\wedge e^J$,
$T^2=T_I\wedge T^I$, and $P_H$ analogously for the Lorentz block.  Direct exterior
algebra gives
\begin{equation}
 P_G=P_H+2\lambda H-\lambda T^2
 =P_H-\lambda N_{\rm NY},\qquad
 N_{\rm NY}=T^2-2H=d(e_I\wedge T^I).
 \label{eq:nieh-yan}
\end{equation}
This is the Nieh--Yan identity \cite{NiehYan}.  Therefore the apparent Holst-type cross
term in the strict full-$G$ Killing invariant is tied to torsion square into a topological
direction.  Combining it with Theorem \ref{thm:invariant-count}, the bulk action space in
the declared class is one-dimensional modulo topological terms.  This does not constitute
a universal Immirzi-parameter no-go \cite{Holst,Immirzi}; $DV$, reduced-$H$ projectors,
fermions, higher powers, and boundaries lie outside the class.

\section{Variational rank and the classical gravity bridge}

Treat the fixed-norm contact reduction as supplied and take compactly supported
variations of the reduced connection and coframe.
The contact direction is not varied as an unconstrained extra field.
In \eqref{eq:explicit-action} the Euler term is locally topological.
Using $\delta R=D^\omega\delta\omega$ and integrating by parts, the
connection variation is proportional to
$\epsilon_{IJKL}e^K\wedge T^L$.
The coframe variation is
\begin{equation}
 \delta_e S_\epsilon
 =4\beta\lambda\int\epsilon_{IJKL}
 \delta e^I\wedge e^J\wedge
 (R^{KL}+\lambda e^K\wedge e^L).
\end{equation}
Here $\beta\lambda\ne0$ is assumed.
For an invertible coframe, variation of \eqref{eq:mm-action} gives the connection equation
\begin{equation}
 \epsilon_{IJKL}e^K\wedge T^L=0.
 \label{eq:torsion-equation}
\end{equation}
The Hodge dual on internal bivectors is invertible, so this is equivalent to
$e^I\wedge T^J-e^J\wedge T^I=0$.
Let $\iota_J$ denote contraction with the dual frame and set
$t=\sum_J\iota_JT^J$.
Contracting the preceding identity and summing over $J$ in four dimensions
gives $T^I=-e^I\wedge t$.
A second contraction gives $t=-3t$, hence $t=0$ and $T^I=0$.
The exact $24\times24$ component map has rank 24 and determinant 16 in the
appendix basis, providing an independent finite certificate.

The coframe equation is
\begin{equation}
 \epsilon_{IJKL}e^J\wedge F_h^{KL}=0,
 \qquad F_h^{IJ}=R^{IJ}+\lambda e^I\wedge e^J.
 \label{eq:coframe-equation}
\end{equation}
For torsion-free data, $F_h$ is an algebraic curvature tensor.
The Hodge dual of this equation is proportional to its Einstein contraction.
In four dimensions, the trace of a vanishing Einstein contraction forces the
scalar curvature of $F_h$ to vanish, and then its Ricci tensor vanishes.
Its unrestricted part is exactly Weyl curvature.
Consequently the linear map on the 20-dimensional residual algebraic
Riemann space has rank ten and kernel dimension ten.
The appendix supplies an exact reference-signature component certificate
and explains why its rank transfers to Lorentz signature.
In the stated internal convention the condition is
$\operatorname{Ric}(R)=-3\lambda g$.
Thus the action imposes Einstein curvature, not $F_h=0$; a constant reversal
of metric sign changes the scalar curvature label and must be carried
separately when assigning physical de Sitter names.

For completeness, let $e_x:T_xM\to P\times_H\mathbb R^{1,3}$ be a pointwise isomorphism
and define
\begin{equation}
 g(X,Y)=\eta(e(X),e(Y)),\qquad
 e(\nabla_XY)=D_X^\omega e(Y).
 \label{eq:solder-connection}
\end{equation}
Then $D^\omega\eta=0$ implies $\nabla g=0$, while
\begin{equation}
 e(T^\nabla)=D^\omega e.
 \label{eq:torsion-lock}
\end{equation}
Full-rank soldering plus \eqref{eq:torsion-equation} therefore makes $\nabla$ the unique
Levi--Civita connection, and the spin curvature becomes the Riemann curvature under the
solder map.  This is the conditional point where an internal Cartan gauge field becomes
metric gravity.

\section{What remains open}

The algebra fixes neither a contact direction nor a physical value of $\ell$.  A
$G$-equivariant selector from no data to $G/H$ cannot exist, because no point of $G/H$ is
fixed by all of $G$.  The contact field must be an order parameter, state, boundary
condition, or other new datum.  A minimal norm potential fixes the radial direction but
does not select its scale dynamically, and the invariant kinetic scalar
$g^{\mu\nu}D_\mu V^AD_\nu V_A=4$ becomes only a volume term once the same $DV$ defines an
invertible metric.

The positive carrier also does not supply Born normalization for noncompact congruences,
the overall coefficient $1/\kappa$, a coherent phase, matter representations and
couplings, a global Spin lift, a Hamiltonian constraint algebra, or a quantum measure.
These are not small corrections to Theorems \ref{thm:qubit-obstruction}--
\ref{thm:invariant-count}; they are the remaining physical theory.

\section{Conclusion}

One qubit positive cone is too small for the simple ten-generator Cartan algebra.  A
four-component spinor is the minimal faithful linear escape, but its gravitational contact
orbit is necessarily signed, and every nonzero Lorentz-fixed Hermitian
selector in this carrier is indefinite even at the positive-cone boundary.
In the narrow local polynomial class stated above, the same
carrier selects a unique contact-linear curvature-square density, while the second
full-Cartan quadratic direction is topological and obeys the Nieh--Yan identity.  Exact
variational ranks then recover torsion-free Einstein--Cartan equations conditional on a
supplied full-rank contact reduction.  The result is a coherent representation-theoretic
and classical-action module; it is not a derivation or quantization of gravity.

\appendix
\section{Explicit finite certificate definitions}
\label{app:certificates}

The companion program \texttt{code/verify.py} constructs every matrix from
the definitions below and writes the compact \texttt{data/results.json}.
It needs Python, SymPy and NumPy, with no historical paths or network access.
No large stored matrix archive is needed for reproduction.

\subsection{The invariant constraint and its exact rank}

Let $p_1,\ldots,p_{10}$ enumerate the pairs $0\leq A<B\leq4$
lexicographically.
Use the basis $e_A$ of the vector representation and the 55 commutative
products $p_ap_b$, $a\leq b$, of its bivector basis.
For each generator $J$, define its bivector action by the Leibniz rule
\begin{equation}
 J(v\wedge w)=(Jv)\wedge w+v\wedge(Jw),
\end{equation}
and its symmetric-square action by
$J(p_ap_b)=(Jp_a)p_b+p_a(Jp_b)$.
The integer action on the 275-dimensional tensor space is
\begin{equation}
 M_J=J\otimes I_{55}+I_5\otimes J_{\mathrm{Sym}^2}.
\end{equation}
Stacking the ten $M_J$ gives an integer $2750\times275$ matrix $M$.
In the stated basis an invariant vector has coefficients proportional to
\begin{equation}
 v_{A,(a,b)}=\epsilon_{A\,p_a^{(1)}p_a^{(2)}p_b^{(1)}p_b^{(2)}} .
 \label{eq:epsilon-null}
\end{equation}
All its nonzero entries involve distinct pairs and receive the same
irrelevant symmetric-product normalization.
The code verifies $Mv=0$ over the integers with $v\ne0$.
It also performs exact modular Gaussian elimination modulo the prime 1009,
obtaining rank 274 for each real form.
A nonzero modular minor implies rational rank at least 274, while
\eqref{eq:epsilon-null} gives rational rank at most 274.
Thus the rational rank is exactly 274.
The invariant coefficient problem is dual to this representation through its
nondegenerate invariant pairing and has the same invariant dimension.
This is an exact argument, independent of singular-value thresholds.

\subsection{Torsion and coframe maps}

Set $e^I=dx^I$ at one point.
Index the torsion columns by $(L,bc)$ with $b<c$, and its components by
$T^L=\frac12 T^L_{bc}dx^b\wedge dx^c$.
The connection-equation rows are $(IJ,abc)$ with $I<J$ and $a<b<c$.
Their entries are
\begin{equation}
 (\mathcal T T)_{IJ,abc}
 =\sum_{K,L}\epsilon_{IJKL}
 \bigl(\delta^K_aT^L_{bc}
      +\delta^K_bT^L_{ca}
      +\delta^K_cT^L_{ab}\bigr).
 \label{eq:torsion-matrix}
\end{equation}
The resulting $24\times24$ integer matrix has exact rank 24 and determinant
16 in these lexicographic conventions.
The determinant is a basis-dependent check, whereas invertibility is the
geometric statement.

For the coframe certificate, begin in a Euclidean orthonormal reference
basis.
An algebraic curvature is a symmetric $6\times6$ matrix indexed by
unordered bivector pairs, giving 21 components before the first Bianchi
identity.
In four dimensions that identity is the single independent condition
\begin{equation}
 R_{01,23}-R_{02,13}+R_{03,12}=0.
\end{equation}
The code obtains its 20-dimensional kernel by exact rational elimination,
not singular-value decomposition.
On this kernel the reference-signature coframe map is
\begin{equation}
 (\mathcal E R)_{I,abc}
 =\sum_{J,K,L}\epsilon_{IJKL}
 \bigl(\delta^J_aR_{KL,bc}
      +\delta^J_bR_{KL,ca}
      +\delta^J_cR_{KL,ab}\bigr).
 \label{eq:coframe-matrix}
\end{equation}
It has rank ten and kernel dimension ten.

Equation~\eqref{eq:coframe-matrix} is explicitly a reference-signature
certificate, rather than an array advertised as Lorentz components with
metric-raising signs omitted.
Complexifying the actual covariant tensor spaces identifies all
nondegenerate metric signatures by invertible complex coframe changes.
The epsilon tensor changes only by a nonzero determinant factor.
Therefore these maps have the same complex rank, and a real matrix has the
same rank over the reals as over its complexification.
The rank result consequently transfers to Lorentz signature.
The identification of the Lorentzian kernel with Weyl curvature is the
separate analytic Einstein-contraction argument in the main text.

\subsection{Clifford and selector checks}

The program checks all Clifford anticommutators, all spin-generator
commutators, every component of the epsilon trace, contact Hermiticity,
contact covariance, and the determinant identity as exact symbolic
equalities.
At unit contact length its spectrum is exactly $\{1,1,-1,-1\}$.
The common Lorentz-fixed matrix space has complex dimension two; the
Hermitian fixed-space classification and its positivity conclusion are
proved in Theorem~\ref{thm:boundary-selector}.

\section{Evidence level and chronology}

The analytic proofs and rational/modular checks establish finite algebraic
claims.
They do not test a global bundle reduction, a global Spin lift, a quantum
measure, a continuum limit or a dynamical selection of the contact field.
The original private records used floating rank and singular-value
calculations for the invariant and variational matrices despite some
exact-certificate wording.
This edition supplies the exact portable reconstruction above rather than
inheriting those labels.

The local representation/action records are dated 13 July 2026.
The present English edition, the exact rank reconstruction, and the explicit
positive-boundary selector corollary were prepared on 29 October 2026.
Private dates do not establish public priority.
The operator reports that earlier work was led by Fable~5, while results from
July 2026 onward were predominantly GPT-produced; DeepSeek and multiple GPT
models also contributed review and preparation.
Human authors and affiliations remain blank.
Human expert review has not occurred and no global novelty certificate is
claimed.
The operator's attribution statement is retained verbatim:
\begin{quote}
I genuinely do not know how authorship, intellectual credit, or ownership
should be attributed for these results.
\end{quote}

\bibliographystyle{unsrt}
\bibliography{references}
\end{document}
